\chapter{Dental Problems}

\section*{Definition and Overview}
Dental problems are conditions that affect the teeth, gums, and mouth, ranging from common issues such as tooth decay and gum disease to toothache, sensitivity, bad breath, and tooth loss. Dental problems cause pain and discomfort, affect eating and speaking, and can have wider effects on health and confidence. Most dental problems are preventable with good oral hygiene, a healthy diet, and regular dental care, and most are highly treatable when detected early. This chapter provides an overview of the most common dental problems, their causes, and what to do about them.

\section*{Epidemiology}
Dental problems are among the most common health conditions worldwide. Tooth decay affects billions of people, and severe gum disease affects around one in five adults globally. In Australia, a majority of adults have had tooth decay, and gum disease affects about one in three. Dental problems are a leading cause of lost work and school time, and they increase with age and with social disadvantage. Access to affordable dental care is a significant factor in who experiences dental disease and its complications.

\section*{Aetiology and Pathophysiology}
\begin{itemize}
    \item \textbf{Tooth decay:} plaque bacteria ferment sugar, producing acid that dissolves enamel and creates cavities
    \item \textbf{Gum disease:} plaque at the gum line causes gingivitis, which can progress to periodontitis with bone loss
    \item \textbf{Toothache:} usually from pulpitis (inflamed pulp) or a dental abscess
    \item \textbf{Sensitivity:} exposed dentine from gum recession, erosion, or wear causes pain on hot, cold, sweet, or sour stimuli
    \item \textbf{Bad breath:} bacteria breaking down food debris in the mouth, often with gum disease or poor hygiene
    \item \textbf{Tooth wear:} erosion from acid, abrasion from brushing, and attrition from grinding
    \item \textbf{Trauma:} chipped, cracked, or knocked-out teeth from injury
    \item \textbf{Other problems:} impacted wisdom teeth, mouth ulcers, oral thrush, and oral cancer
\end{itemize}

\section*{Clinical Features}
\begin{itemize}
    \item Toothache, which may be sharp, throbbing, or constant
    \item Sensitivity to hot, cold, sweet, or sour foods and drinks
    \item Bleeding, swollen, or receding gums
    \item Bad breath or a bad taste
    \item Visible holes, dark spots, chips, or broken teeth
    \item Loose teeth or gaps where teeth have been lost
    \item Mouth ulcers, white patches, or lumps that do not heal
    \item Pain on chewing or difficulty opening the mouth
\end{itemize}

\section*{Differential Diagnosis}
\begin{itemize}
    \item Decay versus staining, fluorosis, or enamel defects
    \item Pulpitis versus dental abscess versus sinus pain causing toothache
    \item Gingivitis versus periodontitis versus other causes of gum swelling
    \item Tooth sensitivity from erosion versus abrasion versus grinding
    \item Mouth ulcers from trauma versus infections, immune conditions, or cancer
    \item Facial pain from dental causes versus sinusitis, jaw joint disorders, or neuralgia
\end{itemize}

\section*{When to Seek Medical Care}
See a dentist for any persistent dental problem, including toothache, bleeding gums, sensitivity, bad breath, or visible changes. Many dental problems are painless in the early stages, so regular check-ups matter. Seek urgent care for severe pain, swelling, or signs of infection.

\textbf{Call 000 (or local emergency services) for:}
\begin{itemize}
    \item Facial or neck swelling with difficulty breathing or swallowing
    \item High fever with dental pain and swelling (possible spreading infection)
    \item A tooth knocked out or a serious mouth injury
    \item Uncontrolled bleeding from the mouth
\end{itemize}

\section*{Diagnosis and Investigations}
\begin{itemize}
    \item \textbf{Dental history:} the nature, timing, and triggers of the problem
    \item \textbf{Clinical examination:} inspection of the teeth, gums, and soft tissues
    \item \textbf{X-rays:} detect decay, abscesses, bone loss, and impacted teeth
    \item \textbf{Vitality testing:} assessment of the pulp of symptomatic teeth
    \item \textbf{Periodontal examination:} measurement of gum pockets
    \item \textbf{Biopsy:} for suspicious lumps or ulcers in the mouth
\end{itemize}

\section*{Management}
\begin{itemize}
    \item \textbf{Decay:} fillings, crowns, root canal treatment, or extraction depending on severity
    \item \textbf{Gum disease:} professional cleaning, scaling and root planing, and improved home care
    \item \textbf{Toothache:} dental treatment of the cause, with pain relief in the meantime
    \item \textbf{Sensitivity:} desensitising toothpaste, fluoride, and treatment of the underlying cause
    \item \textbf{Bad breath:} improved hygiene, tongue cleaning, and treatment of gum disease
    \item \textbf{Trauma:} repair, splinting, or re-implantation of injured teeth
    \item \textbf{Wisdom teeth:} monitoring or extraction for impacted or infected teeth
    \item \textbf{Oral lesions:} treatment of thrush, ulcers, and other conditions, with referral for suspected cancer
\end{itemize}

\section*{Complications}
\begin{itemize}
    \item Spread of dental infection to the face, neck, or bloodstream
    \item Tooth loss and the need for bridges, implants, or dentures
    \item Chronic pain and difficulty eating
    \item Impact on speech, appearance, and self-esteem
    \item Links between gum disease and conditions such as diabetes and heart disease
    \item Poor nutrition in older adults with dental problems
    \item Cost of complex treatment for advanced disease
\end{itemize}

\section*{Prognosis and Outcomes}
Most dental problems are preventable and, when treated early, reversible or fully managed. Tooth decay and gum disease respond well to treatment plus improved prevention, and even advanced disease can be treated to preserve teeth or replace them successfully. Regular dental care and good habits keep the mouth healthy for life. The outlook is excellent for people who attend check-ups and treat problems promptly.

\section*{Prevention and Self-Care}
\begin{itemize}
    \item Brush twice daily with fluoride toothpaste and floss daily
    \item Limit sugary and acidic foods and drinks
    \item Attend regular dental check-ups
    \item Do not smoke, and limit alcohol
    \item Wear a mouthguard for sport
    \item See a dentist promptly for pain, bleeding, or other symptoms
    \item Report any mouth sore, lump, or white patch that does not heal within two weeks
\end{itemize}

\section*{Sources and Further Reading}
The following reputable sources provide further information for healthcare students and patients seeking reliable, current advice about dental problems.

\begin{itemize}
    \item Australian Dental Association. \textit{Dental problems: information for patients.} https://www.ada.org.au/
    \item NHS. \textit{Dental problems: symptoms and treatment.} https://www.nhs.uk/conditions/teeth-problems/
    \item Mayo Clinic. \textit{Oral health: common problems and treatments.} https://www.mayoclinic.org/
    \item World Health Organization. \textit{Oral health fact sheet.} https://www.who.int/
    \item Healthdirect Australia. \textit{Dental and oral health information.} https://www.healthdirect.gov.au/
    \item MedlinePlus. \textit{Tooth and mouth disorders.} https://medlineplus.gov/
\end{itemize}
